\chapter{Background and Theory}
\label{ch:background}

This chapter reviews, without reference to any particular implementation,
the theory that \cref{ch:validation,ch:sta,ch:logical-effort,ch:power-area,ch:optimization,ch:monte-carlo}
implement. Notation follows the assignment specification.

\section{The Circuit as a Timing Graph}
\label{sec:bg-graph}

A flattened combinational netlist is modeled as a directed acyclic graph
$\mathcal{G} = (\mathcal{V}, \mathcal{E})$. Each vertex is either a primary
input, a gate instance, or (implicitly, via a flag on the driving gate) a
primary output; each directed edge represents a net connecting the output of
one vertex to an input pin of another. Because the circuit is purely
combinational, $\mathcal{G}$ must be acyclic --- a directed cycle would
correspond to a signal depending on its own future value, which is not
realizable in combinational logic and must be rejected as a structural
error (\cref{sec:val-cycle}).

Because $\mathcal{G}$ is a DAG, it admits at least one \emph{topological
order}: a linear ordering of vertices such that every edge points from an
earlier vertex to a later one. Any algorithm that must process every gate
only after all of its fan-in gates have already been processed --- as the
forward timing pass in \cref{sec:sta-forward} does --- can be implemented as
a single linear pass over a topological order, in $O(|\mathcal{V}| +
|\mathcal{E}|)$ time, without needing to re-discover dependency order on the
fly.

\section{Load and Delay Modeling}
\label{sec:bg-delay}

\subsection{Output Load Capacitance}

The propagation delay of a logic gate depends on the total capacitance it
must drive at its output. For gate $i$, let $\mathrm{Fanout}(i)$ be the set
of $(j,p)$ pairs such that gate $i$'s output drives input pin $p$ of gate
$j$; let $C_{\mathrm{input},j,p}$ be the input pin capacitance of pin $p$ of
the currently-assigned cell for gate $j$; and let $C_{\mathrm{external},i}$
be any additional load applied directly to gate $i$'s output net (nonzero
only when that net is a primary output, in which case its value is supplied
by the configuration file). The total output load capacitance of gate $i$
is

\begin{equation}
    C_{\mathrm{load},i} \;=\; \sum_{(j,p)\,\in\,\mathrm{Fanout}(i)} C_{\mathrm{input},j,p} \;+\; C_{\mathrm{external},i}.
    \label{eq:cload}
\end{equation}

\subsection{Propagation Delay}

A linear delay model based on effective output resistance and load
capacitance is used. Because the pull-up and pull-down networks of a static
CMOS gate are generally not symmetric, rising and falling output transitions
are modeled with independent parameters: an intrinsic (load-independent)
delay $D_{\mathrm{intrinsic},\mathrm{rise}/\mathrm{fall},i}$ and an
effective output resistance $R_{\mathrm{rise}/\mathrm{fall},i}$, both
supplied per cell by the standard-cell library, together with a
technology-dependent unit-conversion constant $\kappa$:

\begin{align}
    D_{\mathrm{rise},i} &= D_{\mathrm{intrinsic},\mathrm{rise},i} + \kappa\, R_{\mathrm{rise},i}\, C_{\mathrm{load},i}, \label{eq:drise}\\
    D_{\mathrm{fall},i} &= D_{\mathrm{intrinsic},\mathrm{fall},i} + \kappa\, R_{\mathrm{fall},i}\, C_{\mathrm{load},i}. \label{eq:dfall}
\end{align}

Both $D_{\mathrm{rise},i}$ and $D_{\mathrm{fall},i}$ depend only on the
gate's own currently-assigned cell and its total output load
(\cref{eq:cload}); neither depends on which input transition triggered the
output change, nor on any arrival time. This decoupling is what allows load
and delay to be computed for every gate in the circuit in a single pass,
independent of, and prior to, the timing-propagation passes described next.

\section{Unate Timing Arcs}
\label{sec:bg-unate}

A gate's output transition is not always the same polarity as the input
transition that caused it; the relationship is called the pin's
\emph{timing sense} and is characterized per input pin by the standard-cell
library as one of three classes.

\begin{itemize}
    \item \textbf{Positive unate.} The output transition has the same
          polarity as the input transition: an input rise causes an output
          rise, and an input fall causes an output fall. Buffers and
          AND/OR-type gates (an even number of inversions from that input
          to the output) exhibit this behavior.
    \item \textbf{Negative unate.} The output transition has the opposite
          polarity: an input rise causes an output fall, and vice versa.
          Inverters and NAND/NOR-type gates (an odd number of inversions)
          exhibit this behavior.
    \item \textbf{Non-unate.} Either input polarity can, depending on the
          state of the gate's other inputs, cause either output polarity.
          XOR/XNOR-type gates exhibit this behavior, and both possibilities
          must be considered by the timing engine rather than assuming a
          fixed relationship.
\end{itemize}

This classification determines which (input transition, output transition)
pairs are legal timing arcs through a given pin, and both the forward and
backward propagation equations in \cref{sec:bg-sta} must restrict their
maximization/minimization to only the legal pairs for each pin.

\section{Static Timing Analysis Equations}
\label{sec:bg-sta}

\subsection{Forward Propagation --- Arrival Time}
\label{sec:bg-sta-forward}

Let $\mathrm{AT}_i^{s}$ denote the arrival time of transition $s \in
\{\mathrm{rise}, \mathrm{fall}\}$ at the output of gate $i$. For each input
pin $p$ of gate $i$, let $j = \mathrm{driver}(i,p)$ be the gate (or primary
input) driving that pin, and let $\mathcal{T}_{i,p}$ be the set of legal
$(s', s)$ transition pairs for that pin, as determined by its timing sense
(\cref{sec:bg-unate}). Because the propagation delay of gate $i$
(\cref{eq:drise,eq:dfall}) depends only on gate $i$'s own output transition
$s$ and not on which input transition produced it, the arrival time at gate
$i$'s output is

\begin{equation}
    \mathrm{AT}_i^{s} \;=\; \max_{\substack{p \in \mathrm{Inputs}(i) \\ j = \mathrm{driver}(i,p) \\ (s',s)\,\in\,\mathcal{T}_{i,p}}} \left( \mathrm{AT}_j^{s'} + D_i^{s} \right), \qquad s \in \{\mathrm{rise},\mathrm{fall}\}.
    \label{eq:forward-sta}
\end{equation}

Primary inputs are the base case of this recursion: their arrival times are
supplied directly by the configuration file rather than computed. Evaluating
\cref{eq:forward-sta} for every gate, in topological order, yields the
arrival time of every net in the circuit in a single linear pass.

\subsection{Backward Propagation --- Required Time}
\label{sec:bg-sta-backward}

Let $\mathrm{RT}_j^{s'}$ denote the required time of transition $s'$ at the
output of gate $j$: the latest moment that transition may occur without
causing a violation at any downstream primary output. The base case is now
at the primary outputs, whose required times are supplied directly by the
configuration file. For an internal gate $j$, let $\mathrm{Fanout}(j)$ be
the set of $(i,p)$ pairs such that $j$'s output drives pin $p$ of gate $i$.
Because $j$'s output must be ready in time to satisfy \emph{every} path that
depends on it, the required time is the \emph{minimum} over all of $j$'s
fan-out arcs:

\begin{equation}
    \mathrm{RT}_j^{s'} \;=\; \min_{\substack{(i,p)\,\in\,\mathrm{Fanout}(j) \\ (s',s)\,\in\,\mathcal{T}_{i,p}}} \left( \mathrm{RT}_i^{s} - D_i^{s} \right), \qquad s' \in \{\mathrm{rise},\mathrm{fall}\}.
    \label{eq:backward-sta}
\end{equation}

Evaluating \cref{eq:backward-sta} for every gate, in \emph{reverse}
topological order, yields the required time of every net in the circuit.

\subsection{Slack, WNS, and TNS}
\label{sec:bg-sta-slack}

The \emph{slack} of transition $s$ at node $i$ is the margin between when
the signal is required and when it actually arrives:

\begin{equation}
    S_i^{s} = \mathrm{RT}_i^{s} - \mathrm{AT}_i^{s}, \qquad S_i = \min\!\left(S_i^{\mathrm{rise}}, S_i^{\mathrm{fall}}\right).
    \label{eq:slack}
\end{equation}

A negative slack indicates a timing violation at that node. Two
circuit-level summary metrics are defined over the set $\mathcal{O}$ of
primary outputs: the \emph{worst negative slack} (WNS), the single most
severe violation anywhere in the circuit, and the \emph{total negative
slack} (TNS), the sum of every violation's magnitude (zero if the circuit is
fully compliant):

\begin{equation}
    \mathrm{WNS} = \min_{o \in \mathcal{O}} S_o, \qquad
    \mathrm{TNS} = \sum_{o \in \mathcal{O}} \min(0, S_o).
    \label{eq:wns-tns}
\end{equation}

The \emph{critical path} for a given endpoint is the sequence of gates whose
delays, summed along the arcs selected by the $\arg\max$ in
\cref{eq:forward-sta}, produced that endpoint's arrival time; the top-$K$
critical paths are the $K$ (endpoint, transition) pairs with least slack,
together with their reconstructed paths.

\section{Logical Effort}
\label{sec:bg-le}

Logical effort is a technology-independent method for estimating and
minimizing the delay of a chain of static CMOS gates.
It decomposes the normalized delay of a single stage $i$ into a
\emph{parasitic} term $p_i$ (the load-independent delay contributed by the
gate's own internal capacitance) and a \emph{stage effort} term $f_i$, the
product of three unitless quantities:

\begin{itemize}
    \item $g_i$, the \textbf{logical effort} of the input pin the path
          traverses: how much worse that gate is at driving a given load
          than a reference inverter of equivalent drive strength, due
          purely to its logic function (e.g.\ a two-input NAND has $g > 1$
          because its series pull-down transistors must be sized larger for
          equal drive strength).
    \item $h_i$, the \textbf{electrical effort} (fan-out ratio): the ratio
          of the load capacitance the stage must drive, along the path, to
          its own input pin capacitance.
    \item $b_i$, the \textbf{branching effort}: the extra ``tax'' incurred
          when a stage's output also drives capacitance that lies off the
          path under analysis.
\end{itemize}

For a stage with on-path receiver capacitance $C_{\mathrm{onpath},i}$ and
off-path (other fan-out) capacitance $C_{\mathrm{offpath},i}$, and own input
pin capacitance $C_{\mathrm{in},i}$:

\begin{equation}
    b_i = \frac{C_{\mathrm{onpath},i} + C_{\mathrm{offpath},i}}{C_{\mathrm{onpath},i}},
    \qquad
    h_i = \frac{C_{\mathrm{onpath},i}}{C_{\mathrm{in},i}},
    \qquad
    f_i = g_i\, b_i\, h_i,
    \qquad
    d_i = f_i + p_i.
    \label{eq:le-stage}
\end{equation}

For an $N$-stage path with path logical effort $G$, path branching effort
$B$, path electrical effort $H$ (ratio of the path's terminal load to its
first stage's input capacitance), total path effort $F$, and total
parasitic delay $P$:

\begin{equation}
    G = \prod_{i=1}^{N} g_i, \qquad
    B = \prod_{i=1}^{N} b_i, \qquad
    H = \frac{C_{L,\mathrm{path}}}{C_{\mathrm{in},1}}, \qquad
    F = G\,B\,H = \prod_{i=1}^{N} f_i, \qquad
    P = \sum_{i=1}^{N} p_i.
    \label{eq:le-path}
\end{equation}

A classical result of logical effort theory is that, for a fixed number of
stages $N$, path delay is minimized when every stage carries \emph{equal}
effort $\hat{f} = F^{1/N}$; the corresponding normalized and absolute
theoretical minimum delays are

\begin{equation}
    d_{\mathrm{path,min}} = N F^{1/N} + P,
    \qquad
    D_{LE,\mathrm{min}} = \tau\, d_{\mathrm{path,min}},
    \label{eq:le-min-delay}
\end{equation}

where $\tau$ is a technology-dependent time unit supplied by the cell
library. A stage whose actual effort $f_i$ exceeds $\hat{f}$ is, by this
theory, under-sized relative to the load it drives and is a natural
candidate for up-sizing; \cref{ch:optimization} uses exactly this
observation to prioritize which gate to resize first. Inverting
\cref{eq:le-stage} for a target effort $\hat f$ gives, for stage $i$, the
continuous \emph{desired} input capacitance that would make its effort
equal to the path target, which can then be mapped onto the nearest
discretely-sized cell(s) actually available in the library:

\begin{equation}
    C^{*}_{\mathrm{in},i} = \frac{g_i\, C_{\mathrm{total},i}}{\hat f},
    \qquad
    C_{\mathrm{total},i} = C_{\mathrm{onpath},i} + C_{\mathrm{offpath},i}.
    \label{eq:le-target-cap}
\end{equation}

\section{Power and Area Estimation}
\label{sec:bg-power}

Total circuit area is the sum of the area of every currently-assigned cell.
Total power is the sum of a static \emph{leakage} component (supplied
directly per cell by the library, independent of switching activity) and a
switching-activity-dependent \emph{dynamic} component. For gate $i$ with
switching activity factor $\alpha_i$, effective switched capacitance
$C_{\mathrm{sw},i}$ (the sum of its output load and its own internal
capacitance), supply voltage $V_{DD}$, and clock frequency $f$:

\begin{equation}
    A_{\mathrm{total}} = \sum_{i \in \mathcal{G}} A_i, \qquad
    C_{\mathrm{sw},i} = C_{\mathrm{load},i} + C_{\mathrm{internal},i}, \qquad
    P_{\mathrm{dynamic}} = \sum_{i \in \mathcal{G}} \alpha_i\, C_{\mathrm{sw},i}\, V_{DD}^2\, f,
    \label{eq:power-area}
\end{equation}
\begin{equation}
    P_{\mathrm{leakage}} = \sum_{i \in \mathcal{G}} P_{\mathrm{leakage},i}, \qquad
    P_{\mathrm{total}} = P_{\mathrm{dynamic}} + P_{\mathrm{leakage}}.
    \label{eq:power-total}
\end{equation}

Short-circuit power is not modeled separately; the assignment specification
explicitly permits this simplification.

\section{Statistical Timing: Process Variation}
\label{sec:bg-mc}

Manufacturing process variation causes the true delay of a fabricated gate
to differ from its nominal (library-specified) value. Two components are
modeled: a \emph{global} component $Z_g$, shared by every gate in a given
manufactured die (modeling die-to-die and within-die systematic variation),
and a \emph{local} component $Z_i$, drawn independently per gate (modeling
random device-level mismatch). For Monte Carlo sample $k$, with relative
variation strengths $\sigma_g$ and $\sigma_l$ supplied by the configuration
file:

\begin{equation}
    Z_g^{(k)}, Z_i^{(k)} \sim \mathcal{N}(0,1), \qquad
    D_{i,s}^{(k)} = D_{i,s,\mathrm{nom}}\left(1 + \sigma_g Z_g^{(k)} + \sigma_l Z_i^{(k)}\right), \quad s \in \{\mathrm{rise},\mathrm{fall}\}.
    \label{eq:mc-delay}
\end{equation}

Because delay is a strictly positive physical quantity, any sample that
would yield a non-positive delay for any gate is discarded and re-drawn.
Given $M$ valid samples, and letting $I^{(k)}_{\mathrm{violation}}$
indicate whether sample $k$'s worst negative slack is negative, the
empirical violation probability and timing yield are

\begin{equation}
    \hat{P}_{\mathrm{violation}} = \frac{1}{M}\sum_{k=1}^{M} I^{(k)}_{\mathrm{violation}}, \qquad
    \hat{Y}_{\mathrm{timing}} = 1 - \hat{P}_{\mathrm{violation}}.
    \label{eq:mc-yield}
\end{equation}

Full STA (\cref{sec:bg-sta}) is re-run once per Monte Carlo sample against
the perturbed delay set to obtain that sample's WNS, TNS, circuit delay, and
critical path, from which the distributional statistics reported in
\cref{ch:monte-carlo,ch:results} are computed.

